\section{Experimental Setup}
\label{sec:6}

\subsection{Datasets and System Corpus}
\label{sec:6.1}

The evaluation corpus comprises 2{,}812 components across seventeen system architectures --- twelve synthetic topologies that form the inductive cross-validation folds, and five real-world reference systems withheld entirely from training --- as detailed in Table~\ref{tab:4}:

\begin{table}[htbp]
\centering
\small
\caption{Evaluation corpus at a glance. The twelve synthetic topologies are the inductive Leave-One-Scenario-Out folds of Table~\ref{tab:7}; the five real-world systems are withheld from every training fold and used only for zero-shot transfer (Section~\ref{sec:7.4}). Per-scenario entity and edge counts, read from the committed topology files and verified against them in continuous integration, are in Supplementary Section~S13.}
\label{tab:4}
\begin{tabular}{lrrrrrrr}
\toprule
\textbf{Regime} & \textbf{$|V|$} & \textbf{$|V_{\text{app}}|$} & \textbf{Topics} & \textbf{Brokers} & \textbf{Hosts} & \textbf{Libs} & \textbf{$|E|$} \\
\midrule
Synthetic evaluation scenarios (11) & 2{,}387 & 1{,}295 & 588 & 60 & 194 & 250 & 10{,}657 \\
Synthetic case study --- ATM (1)    &    74 &    26 &  27 &  5 &   8 &   8 &   261 \\
\midrule
\textbf{Synthetic subtotal (12 LOSO folds)} & \textbf{2{,}461} & 1{,}321 & 615 & 65 & 202 & 258 & \textbf{10{,}918} \\
\textbf{Real-world subtotal (5 systems)}    & \textbf{351}    &   141 & 120 & 16 &  32 &  42 & \textbf{700} \\
\midrule
\textbf{Total} & \textbf{2{,}812} & 1{,}462 & 735 & 81 & 234 & 300 & \textbf{11{,}618} \\
\bottomrule
\end{tabular}
\end{table}

Here, $|V|$ is the sum of all five entity-type counts per scenario and $|E|$ counts every raw structural relationship instance in the scenario specification --- the native substrate the simulation oracles traverse, across all six relation types --- rather than the derived \texttt{DEPENDS\_ON} projection built for GNN training. We note the convention because the generator's own build log omits \texttt{CONNECTS\_TO} and therefore reports lower per-scenario totals.

The five real-world architectures were transcribed from authentic open-source repositories using dedicated architectural adapters. The synthetic scenarios come from a parameterized topology generator: each is fully defined by a committed configuration specifying a random seed, per-entity-type counts, seven-number summaries for application publish and subscribe fan-out, applications per host, library fan-in and topic payload size, and categorical distributions over the three QoS dimensions. Degree distributions and clustering emerge from these parameters rather than being synthetically forced. Supplementary Section~S5 reports the generative parameters per topology and Section~S13 the resulting per-scenario counts; complete configurations ship in the replication package.

\noindent\textbf{Corpus Composition across Experimental Regimes.} 
The experimental evaluation spans three complementary regimes with precisely bounded corpora:
\begin{enumerate}
    \item \emph{In-Distribution Evaluation (Table~\ref{tab:5}):} Evaluated across all twelve distributed architecture scenarios ($n = 12$) from Table~\ref{tab:4} using stratified 60\% train / 20\% validation / 20\% test node splits over five random seeds. This regime establishes the baseline fitting performance of each predictor when training and testing are drawn from the same underlying architectural distribution, complementing the inductive out-of-distribution evaluation.
    \item \emph{Inductive Leave-One-Scenario-Out (LOSO) Cross-Validation (Table~\ref{tab:7}):} Evaluated across twelve distinct inductive folds totaling 2{,}461 components: the seven core synthetic scenarios, four extended domain topologies (Telecom RAN, Industrial SCADA, Real-Time Gaming, and Logistics Fleet) --- 2{,}387 components between them --- and an Air Traffic Management (ATM) network scenario contributing the remaining 74. In each fold, models are trained on eleven graphs and tested zero-shot on the held-out twelfth graph.
    \item \emph{Real-World Architectural Transfer (Table~\ref{tab:9b}; Supplementary Section~S7):} The five open-source real-world systems (Autoware.universe, Cloud Microservices, Train-Ticket, Home Assistant, EdgeX Foundry) are never used as training folds; they are withheld entirely and used strictly for zero-shot architectural transfer validation.
\end{enumerate}

Not every analysis in Section~\ref{sec:7} runs on the full corpus: the sensitivity sweeps and the detection benchmark predate the four extended domains and operate on smaller cached subsets. Because a reader comparing figures across subsections would otherwise have no way to tell which population a number belongs to, Supplementary Section~S12 tabulates the subset behind each analysis. Comparisons are made only within a row of that table, and never across rows.



\subsubsection{Reproducibility of the Corpus}
\label{sec:6.1.repro}

The benchmark corpus is designed to be fully regenerable rather than merely statically archived. Each dataset is deterministically generated from its configuration file via:
\begin{quote}
\texttt{python cli/generate\_graph.py batch --input-dir data/scenarios --output-dir <dir>}
\end{quote}
A companion manifest records the random seed, entity counts, git commit hash, and a SHA-256 cryptographic digest for each emitted topology. Continuous integration regression tests assert that every committed dataset regenerates \emph{byte-identically} from its configuration and that all disk digests match the manifest. This guarantees that third parties can reproduce the exact graphs used in our experiments, rather than simply sampling from similar distributions.

\subsection{Baselines and Evaluated Predictors}
\label{sec:6.2}

We evaluate four primary predictor configurations drawn from three families. Predictor names state the family and the substrate: an \mbox{\texttt{-N}} suffix marks a model trained on the \emph{native} multigraph, its absence the derived Application--Library flow projection, and a \mbox{\texttt{-QoS}} suffix marks a configuration that consumes declared QoS contracts. \emph{SaG} throughout denotes the framework, never an individual predictor.
\begin{enumerate}
    \item \textbf{Heterogeneous graph learning (typed HGT).} \textbf{HGT-QoS} (proposed): relation-specific Heterogeneous Graph Transformer (Section~\ref{sec:4}) ingesting the complete native multigraph with 16-dimensional continuous-categorical edge features that encode middleware QoS contracts. Its ablation \textbf{HGT}, which masks those QoS dimensions, is reported in Section~\ref{sec:7.3.1}.
    \item \textbf{Homogeneous graph learning (untyped GAT).} \textbf{GAT-N-QoS}: homogeneous Graph Attention Network \cite{velickovic2018graph} trained on the identical native multigraph substrate with per-type input projections, but untyped, single-relation message passing. Its edge channel carries the scalar QoS aggregate $w(e)$ --- dimension $0$ of the same 16-D encoding HGT-QoS consumes --- rather than the per-dimension decomposition; no homogeneous architecture in our suite ingests the full 16-D vector. The HGT-QoS--GAT-N-QoS contrast therefore bounds the \emph{joint} contribution of relational typing and per-dimension QoS encoding. Section~\ref{sec:7.3.1} separates the second factor within the typed architecture, and the corresponding unweighted ablation is \textbf{GAT-N}. The \texttt{-N} suffix denotes the native substrate and is load-bearing: the same homogeneous architecture run on the \texttt{DEPENDS\_ON} projection is reported as \textbf{GAT} / \textbf{GAT-QoS}, and that is the pair Table~\ref{tab:5} carries.
    \item \textbf{QoS-weighted structural baseline (training-free).} \textbf{Topo-QoS}: QoS-weighted topological centrality evaluated on the derived application flow projection.
    \item \textbf{Unweighted structural baseline (training-free).} \textbf{Topo}: structural centrality combining unweighted betweenness centrality and articulation point scoring on the flow projection.
\end{enumerate}
\noindent In addition, the out-of-distribution evaluation (Table~\ref{tab:7}) reports \textbf{RM} ($Q(v)$, the deterministic hierarchical quality attribution model of Section~\ref{sec:5}) as a diagnostic reference baseline. RM is not fitted to rank failure impact; its inclusion demonstrates how much learned relational prediction adds over static structural attribution (Section~\ref{sec:1.2}). Furthermore, deterministic RM scoring drives every sensitivity sweep in Section~\ref{sec:7.3}, where closed-form formulations isolate parameter effects from neural training stochasticity.

\subsubsection{Evaluation Substrates and Parity Guarantee}
\label{sec:6.2.substrate}

Predictors in this study operate on matched substrates within their respective families:

\noindent\textbf{Graph Learning Models (GAT-N-QoS, HGT-QoS):} Under Leave-One-Scenario-Out (Table~\ref{tab:7}) both learned predictors ingest the complete native typed multigraph across all five entity types (recorded by the shared \mbox{\texttt{-N}} suffix). In-distribution (Table~\ref{tab:5}), GAT/GAT-QoS consume the Application--Library \texttt{DEPENDS\_ON} projection, confounding typing with multi-entity visibility. Node type reaches GAT-N-QoS only through its per-type projection, with untyped GATConv across edges, whereas HGT-QoS uses relation-specific HGTConv weights. The edge channel is also unmatched: GAT-N-QoS consumes scalar $w(e)$ while HGT-QoS consumes all 16 dimensions (separated in Section~\ref{sec:7.3.1}). Parameter budget ($434{,}620$ vs.\ $28{,}168$) and directionality remain open confounds (Section~\ref{sec:8.4}).

\noindent\textbf{Training-Free Structural Baselines (Topo, Topo-QoS):} Topological baselines are evaluated on the derived Application--Library \texttt{DEPENDS\_ON} projection (Section~\ref{sec:3.2}), as raw multigraphs route messages through topics/brokers, leaving Application nodes with near-zero betweenness.

\noindent\textbf{Ground-Truth Independence Guarantee:} Crucially, \textbf{no predictor in either group accesses $G_{\text{structural}}$}: that raw topology is strictly reserved for simulation oracles (Section~\ref{sec:4.4}), verified by \path{tests/test_independence_guarantee.py}. Regardless of substrate, all variants are scored on an identical, independently resolved Application node set ($V_{\text{app}}$, Section~\ref{sec:6.3}).

\subsection{Evaluation Metrics and Protocols}
\label{sec:6.3}

\noindent\textbf{Figure accessibility.} All figures use the Okabe--Ito palette with distinct markers and hatchings for monochrome legibility.

\noindent\textbf{Ranking Precision:} Evaluated via Spearman $\rho$ and Kendall $\tau$ against simulated impact $I^*(v)$ from the primary oracle (Section~\ref{sec:4.3}).
\textbf{Critical-Set Identification:} Measured via $F_1@K$, Precision@$K$, and Recall@$K$ for top-$K$ components ($K = \text{round}(0.20 \cdot |V_{\text{app}}|)$), which coincide identically as top-$K$ set overlap.
\textbf{Statistical Significance:} Paired Wilcoxon signed-rank tests \cite{wilcoxon1945individual} ($p < 0.05$) and bootstrap 95\% CIs ($B = 2{,}000$) over folds \cite{efron1993introduction,spearman1904proof}. In the 12-fold LOSO design, power floor is $p = 0.00049$. Applying Holm's step-down correction across ten full-population rank contrasts (Sections~\ref{sec:7.1}--\ref{sec:7.3.1}), three survive: Topo-QoS over Topo ($p = 0.0010$), Topo over RM ($p = 0.0005$), and unweighted typing HGT over GAT-N ($p = 0.0010$); QoS-weighted typing ($p = 0.0122$) and QoS edge ablation ($p = 0.0093$) remain nominally significant with 11/12 directional fold consistency.

\noindent\textbf{Pre-registration.} The primary out-of-distribution contrast (HGT-QoS vs.\ \texttt{Topo-QoS} under LOSO, 5 fixed seeds, fold as unit of analysis) was pre-registered in the replication package before results were obtained. As reported in Section~\ref{sec:7.1}, the margin did not clear statistical significance.

\subsubsection{Evaluation Population and Protocols}

Every predictor within an evaluation table is scored on an identical node population, resolved strictly from scenario topology and ground truth --- the \textbf{Application} set ($V_{\text{app}}$) unless noted. Pooling node types conflates distinct base rates and triggers Simpson's paradox (Section~\ref{sec:7.3}).

\noindent\textbf{In-Distribution Evaluation:} Stratified 60\% train / 20\% val / 20\% test node splits over five seeds $\{42, 123, 456, 789, 2024\}$, redrawing partitions and initializations.
\textbf{Inductive LOSO Cross-Validation:} Models train on eleven scenarios and test zero-shot on the held-out twelfth across all 12 folds under uniform 3-layer depth and inner-split early stopping (Section~\ref{sec:8.4}).
\textbf{Real-World Architectural Transfer:} Synthetic-trained models evaluate zero-shot on five open-source systems without fine-tuning.